% =====================================================================
% RFD -- IEEE Signal Processing Letters scaffold draft
% =====================================================================
% PURPOSE: structural scaffold only. The author rewrites every sentence.
% Do not polish this file; it exists so the section order, claim order
% and evidence placement are fixed before the prose is written.
%
% Venue:   IEEE Signal Processing Letters -- 4 pages of technical
%          content (two-column) + 1 page of references only. Optional
%          supplementary pages for proofs / reproduction.
% Model:   G. Maman, O. Yair, D. Eytan, R. Talmon, "Domain adaptation
%          using Riemannian geometry of SPD matrices," ICASSP 2019.
%          Borrowed moves: (i) intro names a single-mean PT method's
%          shortcoming in one sentence, then proposes the complement;
%          (ii) short "Riemannian geometry of SPD" section before the
%          method; (iii) one experiment section; (iv) org paragraph.
%
% All numbers below come from results/final/ (see % SRC tags).
% % TODO(author) marks anything that must be checked or decided.
% =====================================================================
\documentclass[journal]{IEEEtran}
\usepackage{amsmath,amssymb,amsthm}
\usepackage{graphicx}
\usepackage{booktabs}
\usepackage{cite}

\newcommand{\Exp}{\mathrm{Exp}}
\newcommand{\Log}{\mathrm{Log}}
\newcommand{\PT}{\Gamma}

\begin{document}

\title{Riemannian Factor Dynamics: Factor Models for\\
Covariance Sequences with a Moving Baseline}
% TODO(author): title. Alternative: "When Does a Moving Baseline Matter?
% Factor Models for Covariance Sequences on the Bures--Wasserstein Manifold"

\author{Adarsh Arun% TODO(author): full name, affiliation, contact
}

\maketitle

% ---------------------------------------------------------------------
\begin{abstract}
% MOVE 1: field + object. MOVE 2: parent's single-baseline assumption.
% MOVE 3: what RFD does (moving centre + transport). MOVE 4: the boundary
% result with numbers. MOVE 5: honest real-data result.
Factor models for sequences of covariance matrices that respect their
positive-definite geometry measure every observation against a single
Fr\'echet mean. When the baseline covariance drifts, that drift is either
absorbed into the factors or left unexplained. We propose Riemannian
Factor Dynamics (RFD), which replaces the global mean with an estimated
polygonal centre path and parallel-transports tangent deviations along it
into a common reference frame before extracting factors from lagged
autocovariances. In 576 paired Bures--Wasserstein simulations, RFD reduces
median latent-signal error by 42.5--57.8\% when the drift has components
outside the factor space, and costs about 1\% when the centre is fixed or
its drift lies inside that space. On hourly $20\times 20$ crypto-asset
covariances over 2025, RFD lowers mean forecast loss relative to the
fixed-centre model on four measures, but the paired difference is not
statistically significant.
\end{abstract}

\begin{IEEEkeywords}
Covariance matrices, Riemannian manifolds, Bures--Wasserstein metric,
parallel transport, factor models, time series.
\end{IEEEkeywords}

% ---------------------------------------------------------------------
\section{Introduction}
% P1: covariance sequences are everywhere in SP/finance; they are SPD,
% Euclidean treatment breaks positivity. Cite tangent-space methods.
Sequences of covariance matrices arise in array processing, EEG, and
financial risk. Because covariance matrices are symmetric positive-definite
(SPD), treating their entries as Euclidean coordinates ignores their
geometry and can yield invalid reconstructions. A common remedy maps each
matrix to the tangent space at a Fr\'echet mean and works linearly
there~\cite{pennec2006,maman2019}.

% P2: the parent model, stated fairly. Its assumption (P2).
The Riemannian factor model (RFM)~\cite{huang2026rfm} extends this idea to
time series: deviations from a global Fr\'echet mean are modelled by a
small number of latent factors identified from lagged
autocovariances~\cite{lamyao2012}, with convergence guarantees in high
dimension. Its theory assumes that the Fr\'echet mean of every observation
equals the same point.

% P3: the shortcoming in ONE sentence (Maman-style move), then why it bites.
When the baseline covariance itself drifts, a single reference point must
attribute that drift to the factors. Drift that lies inside the factor
space is absorbed at little cost, but drift outside it either consumes
factor capacity or remains as structured error.

% P4: contributions as a short list.
This letter makes three contributions:
\begin{itemize}
\item a factor model with an estimated moving centre path, whose tangent
deviations are parallel-transported along that path into a common frame;
\item a paired simulation study that locates exactly when the moving centre
helps and when it costs, including an identification boundary;
\item a causal hourly forecasting study on 20 crypto assets against the
fixed-centre model and classical baselines.
\end{itemize}
% TODO(author): decide whether the identification point (bandwidth
% decides the centre/factor split) is a contribution or a limitation.

% P5: organisation paragraph (model paper does this; cut if space tight).
Section~\ref{sec:geom} reviews the geometry, Section~\ref{sec:method}
presents RFD, Section~\ref{sec:exp} reports experiments, and
Section~\ref{sec:conc} concludes.

% ---------------------------------------------------------------------
\section{Geometry of SPD Matrices}\label{sec:geom}
% Keep to ~1/3 column. Define only what Sec. III uses:
% tangent space, metric, Exp/Log, Frechet mean, parallel transport.
Let $\mathcal{S}_{++}^{p}$ denote the $p\times p$ SPD matrices. The tangent
space at any $P$ is the space of symmetric matrices, of dimension
$d=p(p+1)/2$. Under the Bures--Wasserstein (BW) metric~\cite{bhatia2019},
\begin{equation}
\langle U,V\rangle_P=\tfrac12\operatorname{tr}\!\big(\mathcal{L}_P[U]\,V\big),
\quad P\,\mathcal{L}_P[U]+\mathcal{L}_P[U]\,P=U .
\end{equation}
The inner product depends on $P$, so tangent vectors at different base
points cannot be compared directly. The Levi-Civita connection of the
metric defines both geodesics, with maps $\Exp_P$ and $\Log_P$, and the
parallel transport $\PT_{P\to Q}$, a linear isometry between tangent
spaces. On a curved manifold, transport depends on the path taken.
% TODO(author): one sentence on WHY BW (parent's optimal-transport
% argument) or state it is inherited; mention AIRM is supported.

% ---------------------------------------------------------------------
\section{Riemannian Factor Dynamics}\label{sec:method}
% Mirror the README's four steps. One equation for the model, one
% sentence per step, one paragraph on identification.

\subsection{Model}
RFD writes an observation $X_t\in\mathcal{S}_{++}^{p}$ as
\begin{equation}\label{eq:model}
X_t=\Exp_{\mu(t)}\!\big(P_t A f_t+\delta_t\big),
\end{equation}
where $\mu(t)$ is a moving centre, $A\in\mathbb{R}^{d\times r}$ holds the
loadings in a reference tangent frame, $f_t\in\mathbb{R}^{r}$ are dynamic
factors, $P_t$ transports from the reference frame to $\mu(t)$ along the
centre path, and $\delta_t$ is tangent noise without lagged dependence.
RFM is the special case $\mu(t)\equiv\mu$.

\subsection{Estimation}
% Step 1: centre path
\emph{Centre path.} Local Fr\'echet means are estimated with temporal
bandwidth $h$ and joined by geodesic segments into a polygonal path with
$K$ cells.
% Step 2: transport
\emph{Common frame.} Each deviation $\Log_{\hat\mu(t)}(X_t)$ is transported
along the polygon to the reference point, giving coordinates
$y_t\in\mathbb{R}^{d}$ measured with one inner product.
% Step 3: factors
\emph{Factors.} As in~\cite{lamyao2012,huang2026rfm}, the loading space is
the leading eigenspace of $\sum_{k=1}^{k_0}\hat\Sigma_y(k)\hat\Sigma_y(k)^\top$,
where $\hat\Sigma_y(k)$ is the lag-$k$ autocovariance.
% Step 4: reconstruction
\emph{Reconstruction.} Fitted components are transported back along the
path and mapped to the manifold with $\Exp_{\hat\mu(t)}$.
% TODO(author): check the eigen-operator matches py/rfd/model.py exactly.

\subsection{Identification}
% The key honest paragraph. Bandwidth = identifying assumption.
A slowly drifting centre and a highly persistent factor can generate the
same observations. RFD separates them by construction: movement slower
than the bandwidth $h$ is assigned to the centre, faster persistent
movement to the factors. As $h\to\infty$ with $K=1$, RFD reduces to RFM;
as $h\to0$, the centre interpolates the data and no factor structure
remains. Drift inside the loading space is therefore only weakly
identified, as the experiments confirm.

% ---------------------------------------------------------------------
\section{Experiments}\label{sec:exp}

\subsection{Paired Synthetic Study}
% SRC results/final/parent_rfd_bw_parity_adjudication/report.md
% Setup in 2-3 sentences, then Table I, then 2 findings.
We compare RFD with the parent implementation on 576 paired BW draws:
six centre regimes, sample sizes $n\in\{240,512,2048,8192\}$, and 24
replicates each. Both methods receive the same data, the true rank $r=2$
and two lags.
% TODO(author): add p, h, K used in the campaign (config/parent_rfd_bw_parity.yaml).

\begin{table}[t]
\centering
\caption{Median reduction in latent-signal error, RFD vs.\ RFM.}
\label{tab:synth}
\begin{tabular}{lrrr}
\toprule
Regime & $n=240$ & $n=8192$ & RFD wins \\
\midrule
Parent home & $-11.7\%$ & $-1.0\%$ & 0/96 \\
Fixed control & $-11.8\%$ & $-1.0\%$ & 0/96 \\
Aligned drift & $-13.0\%$ & $-1.0\%$ & 0/96 \\
Mixed drift & $+38.7\%$ & $+42.5\%$ & 96/96 \\
Orthogonal drift & $+49.8\%$ & $+57.8\%$ & 96/96 \\
Curved drift & $+45.0\%$ & $+55.8\%$ & 96/96 \\
\bottomrule
\end{tabular}
\end{table}

% Finding 1: the boundary.
RFD wins every paired draw when drift has components outside the factor
space, and its loading-projector error falls by 96--98\% at $n=8192$. In
the three controls it loses, but the penalty shrinks from about 12\% at
$n=240$ to 1\% at $n=8192$.
% Finding 2: identification.
In the aligned regime RFD estimates the centre path 60.5\% more accurately,
yet RFM still reconstructs slightly better, because drift inside the
loading space is absorbed as common movement. Better centre estimation
alone does not establish a better decomposition.
% TODO(author): figure -- results/figures/readme/synthetic_boundary.png
% fits one column; decide figure vs table (not both).

\subsection{Hourly Crypto Covariance Forecasting}
% SRC results/final/crypto_forecast/report.md
We forecast hourly $20\times20$ realised covariances built from one-minute
returns for all 8,760 hours of 2025, with tuning on 2024. RFD and RFM share
a rank-19 VAR(1) score head and refit every four weeks on a trailing 26-week
window; forecasts use only information available at each origin. Baselines
are last observation, EWMA and log-SPD HAR.

% Result, stated at the strength the evidence allows.
RFD lowers mean loss relative to RFM by 2.48\% (squared Frobenius),
13.26\% (QLIKE), 2.18\% (squared BW) and 7.69\% (minimum-variance portfolio
variance). However, the 95\% Newey--West intervals for both primary paired
losses cross zero, so the comparison is a statistical tie. Log-SPD HAR has
the lowest mean Frobenius, BW and portfolio losses, and EWMA the lowest
QLIKE. Median losses rise under three measures, which suggests the mean
gain comes from a minority of large-error hours.
% TODO(author): one table or bar figure here, not both
% (results/figures/readme/crypto_benchmarks.png).
% TODO(author): connect to alignment hypothesis -- real drift may lie
% largely inside the factor space (r/d = 19/210 ~ 9% for random drift).
% Only include if the alignment test is actually run.

% ---------------------------------------------------------------------
\section{Conclusion}\label{sec:conc}
% 3 sentences: what RFD is; the boundary; the limitation + next step.
RFD relaxes the fixed-baseline assumption of Riemannian factor models by
estimating a moving centre and transporting deviations along it. The
moving centre matters when baseline drift leaves the factor space and
costs little otherwise. RFD gives up the parent's asymptotic guarantees,
and on real data the centre/factor split is set by the bandwidth; testing
how much real drift lies outside the factor space is the natural next step.

% ---------------------------------------------------------------------
\bibliographystyle{IEEEtran}
\begin{thebibliography}{9}
\bibitem{huang2026rfm}
S.-C. Huang, R. Chen, and Y. Chen, ``A Riemannian factor model for
manifold-valued time series,'' arXiv:2607.28385, 2026.
\bibitem{lamyao2012}
C. Lam and Q. Yao, ``Factor modeling for high-dimensional time series:
Inference for the number of factors,'' \emph{Ann. Statist.}, vol.~40,
no.~2, pp.~694--726, 2012.
\bibitem{pennec2006}
X. Pennec, P. Fillard, and N. Ayache, ``A Riemannian framework for tensor
computing,'' \emph{Int. J. Comput. Vis.}, vol.~66, no.~1, pp.~41--66, 2006.
\bibitem{bhatia2019}
R. Bhatia, T. Jain, and Y. Lim, ``On the Bures--Wasserstein distance between
positive definite matrices,'' \emph{Expo. Math.}, vol.~37, no.~2,
pp.~165--191, 2019.
\bibitem{maman2019}
G. Maman, O. Yair, D. Eytan, and R. Talmon, ``Domain adaptation using
Riemannian geometry of SPD matrices,'' in \emph{Proc. IEEE ICASSP}, 2019.
% TODO(author): verify every reference against the source before submitting.
\end{thebibliography}

\end{document}
